% ============================================================================
% ARBEITSBLATT: Elektromagnetische Induktion (LP-09b)
% Physik Klasse 9 - Kompakt auf 2 Seiten
% Stand: 01.02.2026
% ============================================================================

\documentclass[11pt,a4paper]{article}

\usepackage[utf8]{inputenc}
\usepackage[T1]{fontenc}
\usepackage[ngerman]{babel}
\usepackage{geometry}
\usepackage{helvet}
\usepackage{enumitem}
\usepackage{fancyhdr}
\usepackage{xcolor}
\usepackage{tcolorbox}
\usepackage{amssymb}
\usepackage{siunitx}
\usepackage{qrcode}
\usepackage{array}
\usepackage{textcomp}
\usepackage[hidelinks]{hyperref}

\sisetup{locale=DE, per-mode=symbol}
\renewcommand{\familydefault}{\sfdefault}

\geometry{left=1.4cm, right=1.4cm, top=1.3cm, bottom=1.0cm}

% Fachfarbe Physik
\definecolor{physik}{HTML}{2c5aa0}
\colorlet{fachfarbe}{physik}
\definecolor{protokollrand}{HTML}{2a7a4b}
\definecolor{merkerand}{HTML}{b35c00}

% Kopfzeile
\pagestyle{fancy}
\fancyhf{}
\fancyhead[L]{\small\textbf{Physik Klasse 9} | Elektromagnetische Induktion}
\fancyhead[R]{\small Arbeitsblatt}
\renewcommand{\headrulewidth}{0.4pt}
\setlength{\headheight}{14pt}

% Boxen
\tcbuselibrary{skins}

\newtcolorbox{merkkasten}[1][]{
    colback=white, colframe=fachfarbe,
    coltitle=black, colbacktitle=white,
    fonttitle=\bfseries\small,
    boxrule=1pt, arc=0pt,
    left=6pt, right=6pt, top=4pt, bottom=4pt,
    #1
}

\newtcolorbox{protokollbox}[1][]{
    colback=white, colframe=protokollrand,
    coltitle=black, colbacktitle=white,
    fonttitle=\bfseries\small,
    boxrule=1pt, arc=0pt,
    left=6pt, right=6pt, top=4pt, bottom=4pt,
    #1
}

\newtcolorbox{merkebox}[1][]{
    colback=white, colframe=merkerand,
    coltitle=black, colbacktitle=white,
    fonttitle=\bfseries\small,
    boxrule=1pt, arc=0pt,
    left=6pt, right=6pt, top=4pt, bottom=4pt,
    #1
}

\newcommand{\aufgabe}[1]{\noindent\textbf{\textcolor{fachfarbe}{#1}}}
\newcommand{\luecke}[1]{\underline{\hspace{#1}}}

\setlength{\parindent}{0pt}
\setlength{\parskip}{0pt}

\begin{document}

% ============================================================================
% TITEL
% ============================================================================
\begin{minipage}[t]{0.80\textwidth}
    \vspace{0pt}
    {\Large\bfseries Elektromagnetische Induktion}\\[0.2em]
    {\small Name: \luecke{4.5cm} \hfill Datum: \luecke{2.5cm}}
\end{minipage}
\hfill
\begin{minipage}[t]{0.15\textwidth}
    \vspace{0pt}
    \raggedleft
    \href{https://devbox42.github.io/sciencesim/physik/kl09/motor-induktion/LP-09b-induktion.html}{\qrcode[height=1.4cm]{https://devbox42.github.io/sciencesim/physik/kl09/motor-induktion/LP-09b-induktion.html}}
\end{minipage}
\vspace{0.1em}
\hrule
\vspace{0.4em}

% ============================================================================
% LEITFRAGE
% ============================================================================
\begin{merkkasten}
\textbf{Leitfrage:} Der Motor macht aus Strom Bewegung. \textbf{Geht das auch umgekehrt?}
\end{merkkasten}

\vspace{0.3em}

% ============================================================================
% KASTEN 1
% ============================================================================
\begin{protokollbox}[title=Kasten 1: Simulationsexperiment -- Protokoll]
\small
\renewcommand{\arraystretch}{1.25}
\begin{tabular}{|>{\raggedright\arraybackslash}p{4.2cm}|p{2.2cm}|p{7.5cm}|}
\hline
\textbf{Experiment} & \textbf{Spannung $U$} & \textbf{Beobachtung} \\
\hline
E1: Magnet langsam durch Spule & & \\[0.35cm]
\hline
E2: Magnet schnell durch Spule & & \\[0.35cm]
\hline
E3: Magnet ruht in der Spule & & \\[0.35cm]
\hline
E4: Bewegungsrichtung umkehren & & \\[0.35cm]
\hline
\end{tabular}

\vspace{0.3em}
\textbf{Schlussfolgerung:} Eine Spannung entsteht nur, wenn \luecke{5.5cm}.
\end{protokollbox}

\vspace{0.3em}

% ============================================================================
% KASTEN 2
% ============================================================================
\begin{protokollbox}[title=Kasten 2: Einflussfaktoren -- Was erhöht die Spannung?]
\small
\renewcommand{\arraystretch}{1.2}
\begin{tabular}{|p{3.2cm}|c|c|c|p{4cm}|}
\hline
\textbf{Variable} & \textbf{Einst. 1} & \textbf{Einst. 2} & \textbf{$U$ größer?} & \textbf{Erkenntnis} \\
\hline
Geschwindigkeit & langsam & schnell & $\square$ Ja ~ $\square$ Nein & \\[0.25cm]
\hline
Windungszahl n & 10 & 50 & $\square$ Ja ~ $\square$ Nein & \\[0.25cm]
\hline
Magnetstärke & normal & stark & $\square$ Ja ~ $\square$ Nein & \\[0.25cm]
\hline
Spulenfläche A & klein & groß & $\square$ Ja ~ $\square$ Nein & \\[0.25cm]
\hline
\end{tabular}

\vspace{0.2em}
\textbf{Fazit:} Die Induktionsspannung wird größer durch:
1. \luecke{3cm} ~ 2. \luecke{3cm} ~ 3. \luecke{3cm} ~ 4. \luecke{3cm}
\end{protokollbox}

\vspace{0.3em}

% ============================================================================
% KASTEN 3
% ============================================================================
\begin{merkebox}[title=Kasten 3: Merksatz -- Elektromagnetische Induktion]
\small
\textbf{Definition:} Bewegt man einen \luecke{2.5cm} relativ zu einer \luecke{2.5cm}, so entsteht eine \luecke{3.5cm}. Diesen Vorgang nennt man \textbf{Induktion}.

\vspace{0.3em}
\textbf{Voraussetzung:} Es muss eine \luecke{2.5cm} geben -- ohne Bewegung keine Spannung!

\vspace{0.3em}
\textbf{Kernidee:} Eine Spannung entsteht, wenn sich die Anzahl der \luecke{2.5cm}, die die Spule durchsetzen, \textbf{ändert}.

\vspace{0.3em}
\textbf{Induktionsgesetz (qualitativ):} Die Induktionsspannung $U_{\text{ind}}$ ist umso größer, je ...
\begin{itemize}[leftmargin=1.2cm, itemsep=0pt, topsep=2pt]
    \item[$\triangleright$] \luecke{2.5cm} die Bewegung ist (= schnellere Änderung).
    \item[$\triangleright$] \luecke{2.5cm} der Magnet ist (= mehr Feldlinien).
    \item[$\triangleright$] \luecke{2.5cm} die Spule hat.
    \item[$\triangleright$] \luecke{2.5cm} die Spulenfläche ist (= mehr Feldlinien durchsetzen).
\end{itemize}
\end{merkebox}

% ============================================================================
% SEITE 2
% ============================================================================
\newpage

% ============================================================================
% KASTEN 4
% ============================================================================
\begin{merkkasten}
\textbf{Kasten 4: Vergleich Motor und Generator}

\vspace{0.2em}
\small
Motor und Generator haben den \textbf{gleichen Aufbau}, aber \textbf{umgekehrte Funktionen}.

\vspace{0.3em}
\renewcommand{\arraystretch}{1.3}
\begin{tabular}{|p{3.5cm}|p{4.8cm}|p{5cm}|}
\hline
& \textbf{Motor} & \textbf{Generator} \\
\hline
Energieumwandlung & Elektr. Energie $\rightarrow$ \luecke{1.8cm} & \luecke{1.8cm} $\rightarrow$ Elektr. Energie \\
\hline
Strom & fließt \luecke{1.5cm} & kommt \luecke{1.5cm} \\
\hline
Funktionsweise & Strom rein $\rightarrow$ Bewegung & \luecke{3.5cm} \\
\hline
Beispiel & E-Bike, Ventilator & \luecke{2cm} \\
\hline
\end{tabular}

\vspace{0.3em}
\textbf{Merksatz:} Motor und Generator sind „\luecke{2.5cm}" -- gleicher Aufbau, umgekehrte Funktion!
\end{merkkasten}

\vspace{0.4em}

% ============================================================================
% ÜBUNGEN
% ============================================================================
\aufgabe{Übung 1: Anwendungen der Induktion}

\vspace{0.2em}
{\small
\renewcommand{\arraystretch}{1.2}
\begin{tabular}{|p{3cm}|p{11cm}|}
\hline
\textbf{Anwendung} & \textbf{Wie funktioniert die Induktion hier?} \\
\hline
a) Fahrraddynamo & \\[0.6cm]
\hline
b) Wasserkraftwerk & \\[0.6cm]
\hline
c) Windkraftanlage & \\[0.6cm]
\hline
\end{tabular}
}

\vspace{0.4em}

\aufgabe{Übung 2: Richtig oder falsch?}

\vspace{0.2em}
{\small
\begin{tabular}{|p{12cm}|c|c|}
\hline
\textbf{Aussage} & \textbf{R} & \textbf{F} \\
\hline
a) Ein ruhender Magnet in einer Spule erzeugt eine Spannung. & $\square$ & $\square$ \\
\hline
b) Je schneller die Bewegung, desto größer die Induktionsspannung. & $\square$ & $\square$ \\
\hline
c) Der Generator wandelt Bewegungsenergie in elektrische Energie um. & $\square$ & $\square$ \\
\hline
d) Im Kraftwerk werden Motoren zur Stromerzeugung verwendet. & $\square$ & $\square$ \\
\hline
e) Richtungswechsel der Bewegung ändert das Vorzeichen der Spannung. & $\square$ & $\square$ \\
\hline
\end{tabular}
}

\vspace{0.4em}

\aufgabe{Übung 3: Transferaufgabe}

\vspace{0.2em}
{\small Eine Schülerin schüttelt eine Faraday-Taschenlampe (ohne Batterien) und die Lampe leuchtet.}

\vspace{0.2em}
a) Erkläre, wie die Lampe funktioniert. Nutze den Begriff \textit{Induktion}.

\vspace{1.1cm}

b) Warum leuchtet die Lampe heller, wenn man schneller schüttelt?

\vspace{1.1cm}

% ============================================================================
% SELBSTCHECK
% ============================================================================
\begin{merkkasten}
\textbf{Selbstcheck: Das kann ich jetzt!} \hfill {\scriptsize \textcircled{\scriptsize 1} = unsicher \quad \textcircled{\scriptsize 2} = geht so \quad \textcircled{\scriptsize 3} = sicher}

\vspace{0.2em}
\small
\renewcommand{\arraystretch}{1.1}
\begin{tabular}{p{10.5cm}ccc}
Ich kann erklären, was Induktion ist. & $\square$ & $\square$ & $\square$ \\
Ich weiß, wann eine Induktionsspannung entsteht. & $\square$ & $\square$ & $\square$ \\
Ich kann die vier Faktoren nennen, die $U$ beeinflussen. & $\square$ & $\square$ & $\square$ \\
Ich kann Motor und Generator vergleichen. & $\square$ & $\square$ & $\square$ \\
\end{tabular}
\end{merkkasten}

\end{document}
